% Transcription — fonds Alexandre Grothendieck, shelfmark 104, batch 1 (pages 1-16).
% Established from the facsimile archives/batches/104/batch-01.pdf (a mirror of
% https://grothendieck.umontpellier.fr/104.pdf), per the skill
% transcribe-grothendieck. The folder is sixteen pages and this is its only
% batch.
%
% All sixteen leaves are transcribed: every one of them carries his hand and
% every one of them carries mathematics. There is no unrelated verso anywhere
% in the folder, so there is no gap in the page numbers.
%
% The folder splits cleanly in two, and the split is physical as well as
% intellectual. Pages 1 to 5 are five loose landscape sheets — trial fields,
% drawings and formulas thrown across the width of the leaf, written on
% whatever paper was at hand: page 2 is the back of an envelope addressed to
% him, pages 3 and 5 are the backs of typed pages. Pages 6 to 16 are the run
% proper, in ink, in portrait, under his own pagination 1 to 6. That
% pagination stands at the head of pages 6, 8, 10, 12, 14 and 16 only — the
% rectos — and pages 7, 9, 11, 13 and 15 are the versos of those same five
% leaves and carry no number; the text runs continuously across each pair, so
% the run is complete and unbroken from page 6 to page 16.
%
% The inventory's title, « [Champs (stacks) 1] », is bracketed and is
% therefore the archivists'. Nothing in the folder carries a title of his, and
% the two section headings used here are the edition's and are marked as such.
% What the folder is actually about is narrower and more definite than
% « champs » : it is the combinatorics of a category of models M and of the
% ordered sets of subobjects it generates — the apparatus of Pursuing Stacks,
% written out from first principles.
%
% The run of pages 6 to 16 is one construction and it has a single object: to
% replace a category M of « modèles » by a purely order-theoretic datum, and
% to recover M from it. Page 6 fixes the two hypotheses — every arrow of M is
% mono, every endomorphism of an object of M is the identity — and attaches to
% each D the ordered set D-tilde of its subobjects, so that Hom(D',D) injects
% into D-tilde with image the subobjects isomorphic to D'. Page 7 replaces M
% by the family (I_n) indexed by the ordered set N of isomorphism classes,
% with a strictly increasing « application type » tau_n : I_n --> N whose
% image is N_{<=n} and whose fibre over n is the greatest element d_n, and
% turns Hom(D_{n'}, D_n) into a map phi_{n,n'} from that fibre into
% Hom(I_{n'}, I_n). Page 8 adds the two facts that make the translation exact
% — u_xi carries I_{n'} isomorphically onto (I_n)_{<=xi}, and u_{d_n} is the
% identity — states that N, the I_n, the tau_n and the phi_{n,n'} reconstitute
% M up to equivalence, writes the composition law u_xi u_y = u_zeta that the
% admissible homs must satisfy, and defines from those data the category M_0
% whose objects are N and whose arrows are the admissible homs. Page 9 proves
% that M --> M_0 is an equivalence, indeed an isomorphism when M is rigid, and
% sets out the whole datum inside one large brace. Page 10 turns the
% definition round — such data define a category satisfying (i) and (ii) — and
% gives three examples: the finite totally ordered sets, with I_n the non-empty
% subsets of Delta_n; the pinned cubes, with Card I_n = 3^n; and the
% hemispheres, with I_n = (Delta_n x {+-1}) disjoint {d_n}. Page 11 asks
% whether the I_n^* are regular spherical polyhedra and whether other such
% sequences exist, with a marginal aside that it must be in Coxeter's books.
% Pages 12 to 16 leave the models and treat the objects built out of them: the
% subcategory of M-hat whose objects are covered by their model subobjects,
% the « M-épinglage » of the ordered set K of those subobjects, the two axioms
% (1) and (2) it satisfies, the strictly increasing maps compatible with it,
% the closed subsets and the induced pinning, the finiteness condition on
% maximal elements, and finally the stability of the whole under sums and
% amalgamated sums, closing on an N.B. that finite I_n make the objects of
% finite type correspond to the finite pinned ordered sets.
%
% Pages 1 to 5 are the same construction seen earlier and from further off,
% and they are transcribed at their rank because they are mathematics, not
% because they are tidy. Page 1 is a row of sketched cell complexes and one
% line placing (Cat) between the semi-simplicial sets and the ordered sets.
% Pages 2 and 3 are about the nerve: a two-truncated globular system F_0 <= F_1
% <= F_2 with its s_i, t_i, sigma_i, mu_i^j and the three coherence letters
% a, l, r, and the formula F_n = Fl C_1 x ... x Fl C_n x Ob C_{n+1}. Page 4 is
% a page of labelled conditions a), b), b'), c), c') grouped under E_0, E_1,
% E_2, on when a map of such systems is an equivalence, with pi_0(f) and
% pi_1(f) injective in the margin. Page 5 is the draft of pages 6 to 9 with
% the ordered set called I instead of N, and it carries the cube count
% 1 + n.2 + n(n-1)/2 . 2^2 + ... + 2^n = (1+2)^n that page 10's margin will
% reduce to Card I_n = 3^n.
%
% Notation, fixed once for the whole batch rather than page by page. He draws
% two capitals with a doubled stroke and they must not be confused: a barred N
% for the ordered set of isomorphism classes of objects of M, set here
% \mathbb{N}, and the ordinary N of the natural numbers, set here \mathbf{N} —
% page 10's N.B. says « dans les cas que j'ai en tête, \mathbb{N} = \mathbf{N} »,
% which is the whole point of distinguishing them. He writes I_n with a tall
% serifed capital I, barred top and bottom, and that glyph is what settles the
% right-hand side of « l'ens. ordonné \widetilde{D}_n = I_n » at the foot of
% page 6. « homs. » is his own abbreviation and is kept: page 8 writes « ens.
% des homs. admissibles I_{n'} --> I_n ». « st.b.l.té » on page 8 is
% « stabilité ». « t.f. » on page 16 is « type fini ». « Ob » and « Fl » on
% pages 2 and 3 are objects and arrows. The B of page 12 is drawn with a
% doubled stroke and is set \mathbb{B}, as the typed page under page 5 draws
% it.
%
% Readings flagged and not settled: the verb after « De plus, on se » on page 7,
% read « rappelle » on the sense and not on the ink; the adverb after « Notons
% que » on the same page, read « déjà »; the two words after the dash on page 9,
% read « peut donc »; the noun after « les flèches les » on page 13, read
% « homs. » on the strength of page 8; the last word of page 15, read
% « mêmes »; and the word after « ens. des » on page 5, read « classes ».
% Page 4 is the least legible leaf of the folder and its right-hand column is
% given only where the strokes carry.
%
% Two variances of substance are flagged where they occur and not repaired.
% The triangle at the foot of page 13 writes its right-hand vertex with a K
% although f goes from K to L and page 14 writes the same object L_{<=y}. And
% the marginal count « Card I_n = 2n+1 » beside Ex. 3 on page 10 does not
% agree with the I_n = (Delta_n x {+-1}) disjoint {d_n} written in the body,
% which has 2n+3 elements; the neighbouring marginal note writes
% Delta_{n-1} x {+-1}, which does give 2n+1.
%
% Pass: Opus 5 (claude-opus-5), 2026-09-13 — first pass, unchecked against the pages by a human.
\documentclass[11pt,a4paper]{article}
% Le préambule commun à toutes les transcriptions du fonds Grothendieck.
%
% Il tient en une idée : **l'incertitude se compose, elle ne se résout pas.**
% Une transcription qui lisse un mot douteux a détruit la seule chose pour
% laquelle on la faisait — un lecteur doit pouvoir savoir, à chaque endroit,
% ce qui était sur le feuillet et ce qui a été ajouté. Les six macros qui
% suivent sont donc l'appareil critique, et non de la décoration.
%
% Les mêmes commandes sont reconnues par `scripts/render.mjs`, qui produit la
% vue de lecture du site. Une macro ajoutée ici doit l'être aussi là-bas,
% faute de quoi le rendu échouera bruyamment — ce qui est voulu : un rendu
% silencieusement faux est pire qu'un rendu absent.

\ProvidesPackage{grothendieck}[2026/08/08 Transcriptions du fonds Grothendieck]

\usepackage{amsmath,amssymb,amsthm}
\usepackage{mathrsfs}
\usepackage{stmaryrd}
% Les diagrammes commutatifs sont partout dans ce fonds : c'est la langue
% même de la géométrie algébrique de Grothendieck, et une transcription qui
% les rendrait en prose ne transcrirait rien.
\usepackage{tikz-cd}
% Français par défaut — c'est la langue des feuillets — mais l'anglais est
% chargé aussi : la traduction et le résumé sont des documents anglais, et la
% typographie française (espaces avant les deux-points) y serait une faute.
% Ces documents basculent par \selectlanguage{english} en tête de corps.
\usepackage[english,french]{babel}
\usepackage{fontspec}
\usepackage{geometry}
\geometry{a4paper,margin=25mm,marginparwidth=28mm}
\usepackage{marginnote}
\usepackage{xcolor}
% ulem, not soul. `soul`'s \st and \ul are built on a character-by-character
% scan that XeLaTeX's font machinery breaks: the document fails with an
% undefined control sequence rather than degrading. ulem does the same two jobs
% and is Unicode-engine safe. `normalem` stops it hijacking \emph, which the
% transcriptions use for Grothendieck's own emphasis.
\usepackage[normalem]{ulem}
\usepackage{hyperref}
\hypersetup{colorlinks=true,linkcolor=black,urlcolor=black,pdfborder={0 0 0}}

% --- Métadonnées du lot -----------------------------------------------------
% Reprises telles quelles par le script de rendu, qui les affiche en tête de
% la vue de lecture. Elles fixent ce que le fichier prétend couvrir : sans
% elles, une transcription flotte sans point d'attache dans un fonds de seize
% mille feuillets.

\newcommand{\folder}[1]{\gdef\grothFolder{#1}}
\newcommand{\batch}[1]{\gdef\grothBatch{#1}}
\newcommand{\leaves}[2]{\gdef\grothFirst{#1}\gdef\grothLast{#2}}
\newcommand{\foldertitle}[1]{\gdef\grothFoldertitle{#1}}
\gdef\grothFolder{?}\gdef\grothBatch{?}\gdef\grothFirst{?}\gdef\grothLast{?}\gdef\grothFoldertitle{}

% --- Le résumé --------------------------------------------------------------
%
% Il ouvre la lecture modernisée, sous le titre « Résumé », et il est composé
% en petit. Ce n'est pas un ornement : c'est le seul passage du document qui
% ne soit pas des mathématiques, et un lecteur doit pouvoir le reconnaître
% avant de l'avoir lu — puis le sauter s'il connaît déjà le sujet, ou s'y
% arrêter s'il ne le connaît pas. Le filet qui le suit marque la reprise du
% corps.

\newenvironment{resume}
  {\small\setlength{\parskip}{0.45em}}
  {\par\medskip\hrule\medskip}

% --- Mots-clés ---------------------------------------------------------------
%
% En anglais, en clôture du résumé de la lecture modernisée : le vocabulaire
% moderne sous lequel chercher ce que les feuillets construisent. Le manifeste
% du site (`npm run manifest`) les extrait pour étiqueter la cote entière —
% cette ligne est la source unique des tags, il n'y a pas de fichier à côté.
\newcommand{\keywords}[1]{\par\smallskip\noindent
  {\footnotesize\textsc{Keywords} --- \textit{#1}}\par}

% --- L'appareil critique ----------------------------------------------------

% Le numéro de feuillet, en marge. C'est l'ancre qui permet de régler un
% désaccord sur une formule en regardant *un* feuillet plutôt qu'une chemise
% de six cents. Le site s'en sert aussi pour tourner les pages du fac-similé
% quand on fait défiler la transcription.
\newcommand{\leaf}[1]{%
  \marginnote{\footnotesize\color{gray!70}\textsf{#1}}%
  \label{leaf:#1}\ignorespaces}

% Illisible : on ne devine pas. Un mot inventé qui se lit comme les autres est
% le pire résultat possible.
\newcommand{\ill}{\textcolor{red!60!black}{[\,\ldots\,]}}

% Lecture douteuse : le mot est donné, et signalé comme incertain.
\newcommand{\uncertain}[1]{\uline{#1}}

% Ajout de l'éditeur — une lettre manquante, un mot sous-entendu.
\newcommand{\add}[1]{\textcolor{blue!50!black}{[#1]}}

% Biffé par Grothendieck lui-même. Ce qu'il a rayé fait partie du document :
% une rature dit souvent où la pensée a changé de direction.
\newcommand{\struck}[1]{\sout{#1}}

% Note du transcripteur, sur l'état matériel du feuillet.
\newcommand{\note}[1]{\par\smallskip{\small\color{gray!60!black}\textit{[#1]}}\par\smallskip}

% Note marginale de Grothendieck, distincte du corps du texte.
\newcommand{\marginal}[1]{\par\smallskip\noindent\hspace{1em}%
  {\small\color{gray!60!black}#1}\par\smallskip}

% --- Titre ------------------------------------------------------------------

\AtBeginDocument{%
  \begin{center}
    {\large\bfseries Fonds Alexandre Grothendieck --- cote n\textsuperscript{o}~\grothFolder}\\[2pt]
    {\itshape\grothFoldertitle}\\[4pt]
    {\small feuillets \grothFirst--\grothLast\ (lot \grothBatch)}\\[2pt]
    {\footnotesize\color{gray!60!black}Transcription établie d'après le fac-similé numérisé
     par l'Université de Montpellier.\\
     Les lectures incertaines sont soulignées, les illisibles notées [\,\ldots\,],
     les ajouts entre crochets.}
  \end{center}
  \vspace{1em}}

\endinput


\folder{104}
\batch{1}
\pages{1}{16}
\dating{s.d.}
\watermark{Édition de démonstration}
\foldertitle{[Champs (stacks) 1] : copies de notes manuscrites (s.d.)}

\begin{document}

\section{Feuillets d'essai}

\note{ce titre est de l'édition. Les pages 1 à 5 sont cinq feuillets isolés,
écrits en travers, le feuillet tourné d'un quart de tour, sur du papier de
rencontre ; aucun ne porte de titre. Ce sont des champs d'essai pour la
construction que les pages 6 à 16 écrivent au propre, et c'est à ce titre
qu'ils sont transcrits. L'inventaire donne au dossier le titre
« [Champs (stacks) 1] », entre crochets, donc des archivistes : rien dans le
dossier ne porte de titre de sa main}

\page{1}

\note{le feuillet porte, dans sa moitié supérieure, une rangée de figures :
des segments portant deux, trois puis quatre marques de subdivision ; des
formes en amande hachurées, isolées ou accolées deux à deux ; de petits
carrés divisés en quatre ; et un huit, deux cercles tangents en un point. Ce
sont des complexes cellulaires esquissés, et ils ne sont pas redessinés ici}

$(\text{Ss sets}) \supset (\text{Cat}) \longleftarrow (\text{Préord}) \supset
(\text{Ord}) \supset (\text{Maquettes})$

\note{le dernier terme est écrit au-dessus de \struck{Triangulations}. Les
trois signes lus ici $\supset$ sont tracés d'une seule courbe ; la flèche qui
entre dans (Cat) est la seule du trait à porter une pointe}

\[
F_x \longrightarrow F_y
\qquad\qquad
x \leq y
\]

\page{2}

\note{le feuillet est le dos d'une enveloppe à son adresse ; un tapuscrit
transparaît par endroits. Sa moitié droite est occupée par une rangée de
figures — segments subdivisés et fléchés, amandes hachurées accolées, carrés
divisés en bandes ou en damiers, un huit — qui ne sont pas redessinées ici}

\[
F_0 \overset{s_1}{\underset{t_1}{\leftleftarrows}} F_1
\overset{s_2}{\underset{t_2}{\leftleftarrows}} F_2
\]

\[
F_0 \overset{\sigma_1}{\longrightarrow} F_1 \overset{\sigma_2}{\longrightarrow} F_2
\]

\[
\mu_1^1 \qquad \mu_1^2 \qquad \mu_2^2
\]

\[
a \qquad \ell \qquad r
\]

\note{les trois lettres $a$, $\ell$, $r$ sont écrites l'une sous l'autre, sans
autre indication}

\[
I_{A\cdot S} \longrightarrow S
\]

\[
\Delta_{n_1} \times \Delta
\qquad\qquad
\Delta_{n_1} \times \Delta_{n_2} \times \cdots \times \Delta_{n_n}
\qquad\qquad
\Delta_n \times \Delta_{n'}
\]

\begin{align*}
F_0 &= \mathrm{Ob}\,\Delta_{n'} \\
F_1 &= \struck{\ill{}}\ \mathrm{Fl}(\Delta_n) \times \mathrm{Ob}\,\Delta_{n'} \\
F_2 &= \mathrm{Fl}(\Delta_n) \times \mathrm{Fl}(\Delta_{n'})
\end{align*}

\note{la première ligne porte $\Delta_{n'}$ là où la suite des formules
demanderait $\Delta_n$ ; le mot rayé de la deuxième ligne est repris deux fois
et n'est pas lisible}

\page{3}

\note{le feuillet est le dos d'une page dactylographiée, dont le texte
transparaît en gris sur toute la surface ; il n'est pas transcrit. La moitié
droite porte quatre figures : un prisme hachuré posé sur un quadrillage, un
carré, un cylindre fermé en bas par une amande hachurée, et un carré divisé
en quatre}

\[
\begin{array}{l}
\mathrm{Ob}\,C_1 \\
(\mathrm{Fl}\,C_1 \times \mathrm{Ob}\,C_2) \\
\mathrm{Fl}\,C_2 \times \mathrm{Fl}(C_2) \times \mathrm{Ob}\,C_3
\end{array}
\qquad\longrightarrow\qquad \mathrm{Fl}_0(C_1)
\]

\note{les trois lignes sont accolées à gauche par deux flèches montantes ; la
flèche vers $\mathrm{Fl}_0(C_1)$ part de la première. Entre la deuxième et la
troisième est rayé \struck{$\mathrm{Fl}_0(C_1) \times \mathrm{Ob}\,C_2$}. Le
premier indice de la troisième ligne est surchargé et se lit 2 ; la suite des
formules demande $C_1$}

\[
F_3 = \mathrm{Fl}\,C_1 \times \mathrm{Fl}\,C_2 \times \mathrm{Fl}\,C_3 \times \mathrm{Ob}\,C_4
\]

\[
F_n = \mathrm{Fl}\,C_1 \times \mathrm{Fl}\,C_2 \times \cdots \times \mathrm{Fl}\,C_n \times \mathrm{Ob}\,C_q
\]

\note{le dernier indice est écrit $q$ ; la suite des formules demande
$C_{n+1}$}

\[
O_1 \subset F_1, \quad \struck{B_2}\ O_2 \subset F_2, \quad \cdots \quad O_n \subset F_n
\qquad\qquad \dot{F} \subset F
\]

\[
F_1 \times \cdots \times F_n \times (?)
\qquad\qquad
(F \times I) \sqcup_{\dot{B} \times I} \dot{F}
\]

\[
(O_1, F_1) * (O_2, F_2) \cdots * (O_3, F_{n+1})
\]

\note{les indices de la dernière ligne sont écrits ainsi ; le premier facteur
de la fin se lit $O_3$ là où la série demanderait $O_n$}

\page{4}

\note{c'est le feuillet le moins lisible du dossier : une page d'essai écrite
serré sur toute la largeur, au-dessus d'un tapuscrit qui transparaît. Les
conditions étiquetées portent, la prose qui les relie ne porte pas}

\[
\mathbb{Z}[F_{n-1}] \longleftarrow \mathbb{Z}[F_{n-1}] \overset{t_n - s_n}{\longleftarrow}
\mathbb{Z}[F_n] \longleftarrow \mathbb{Z}[F_{n+1}]
\]

\note{les deux premiers termes portent le même indice sur le feuillet}

\[
D_n \longrightarrow X \qquad\qquad D_n^{\ill{}}
\qquad\qquad \xi \in F_n
\]

\[
t_{n-1}\bigl( (t_n y) - s_n(y) \bigr) - s_{n-1}\bigl( t_n y - s_n y \bigr)
\qquad (= 0)
\]

\note{le zéro est écrit sous une accolade qui coiffe toute l'expression}

$\mathrm{E}_0$ \quad \textbf{a)} $\forall\, s' \in F_0'$, $\exists\, s \in F_0$,
$u \in F_1'$, $f(u) : \ill{} \longrightarrow s'$

\textbf{b)} $\forall\, s, t \in F_0$, $u' : f(s) \longrightarrow f(t)$
\quad ($u' \in F_1'$), \quad $\exists\, u \in F_1$, $u : s \longrightarrow t$
--- $\pi_0(f)$ inj.

$\mathrm{E}_1$ \quad \textbf{b')} $\forall\, s, t \in F_0$,
$u' : f(s) \longrightarrow f(t)$ \quad ($u' \in F_1'$), \quad
$\exists\, u \in F_1$, $u : s \longrightarrow t$, \quad $f(u) \ill{} u'$
\quad ($\in F_2$)

\textbf{c)} \struck{$\forall\, s, t \in F_0$,} $u, v \in F_1$,
$s_1(u) = s_1(v)$, $t_1(u) = t_1(v)$, et $w' : f(u) \longrightarrow f(v)$
\quad ($\in F_2'$), \quad $\exists\, w : \ill{}$ \quad ($\in F_2$)
--- $\pi_1(f)$ inj.

$\mathrm{E}_2$ \quad \textbf{c')} $\forall\, (u,v) \in F_1$,
$s_1(u) = s_1(v)$, $t_1(u) = t_1(v)$, et $w' : f(u) \longrightarrow f(v)$
\quad ($\in F_2'$), \quad $\exists\, w : u \longrightarrow v$ \quad ($\in F_2$)
\struck{\ill{}} tel que $w \sim f(w)$

\note{la dernière ligne écrit $w \sim f(w)$ ; le sens demande $w' \sim f(w)$}

\marginal{($C_i$) $\forall\, w' \in F_i'$, \struck{\ill{}} $\exists\, w \in F_i$
tel que $f(w) \sim w'$. De plus si $i \geq 1$, on peut \ill{} se donner
d'avance \ill{} $u = s(w)$, $t = t(w)$, \ill{} aux seules conditions \ill{}
tels que $f(u) = u'$, \ill{} ; et de plus, si $i - 1 \geq 1$, i.e. $i \geq 2$,
après $s(u) = s(v)$, $t(u) = t(v)$. (N.B. \ill{} pour que $f(u) = u'$.)}

\note{cette colonne occupe toute la marge droite du feuillet et n'est lisible
que par fragments}

\page{5}

\note{le feuillet est le dos d'une page de son propre tapuscrit anglais, dont
une phrase reste lisible en travers de la marge gauche : elle pose
$\underline{B} = (\text{Gr-stacks})$, dit que les limites inductives s'y
calculent composante par composante, et précise que $\underline{B}(0)$ désigne
la sous-catégorie pleine des objets de base $\underline{D}_n$, $n \in
\underline{N}$. C'est cette page-là qui nomme les $\widetilde{D}_n$ dont sa
main se sert ici, et c'est pourquoi elle est signalée}

\[
\widetilde{D} \qquad\qquad \mathbb{N} \ni n \longmapsto \widetilde{D}_n
\qquad\qquad \widetilde{D}_{n'} \rightrightarrows \widetilde{D}_n
\]

\[
d_n \in \widetilde{D}_n \longrightarrow \struck{\mathbb{N}}\ I
\qquad\qquad d_n \longmapsto n
\]

\note{à droite de ce trait : « strictement croissante, image
\struck{$\mathbb{N}$} $I_{\leq n}$ »}

\[
\widetilde{D}_n \longrightarrow n \in I
\]

\[
1 + n\cdot 2^1 + \frac{n(n-1)}{2}\,2^2 + \cdots + 2^n = (1+2)^n
\]

$I$ un ordonné

$n \longmapsto \widetilde{D}_n$ famille d'ens. ordonnés

$\widetilde{D}_n \overset{\tau_n}{\longrightarrow} I$ \quad str.\ croiss.,
image $I_{\leq n}$

Si $n' \leq n$, \struck{soit l'ens.\ \ill{}} soit
$\Phi_{n',n} \subset \widetilde{D}_n$, $\ = \{x \in \widetilde{D}_n \mid
\tau_n(x) = n'\}$

\struck{Soit $x \in \Phi_{n'}$} Considérons l'ens.\ \struck{\ill{}}

\[
F_0 \leftleftarrows F_1 \leftleftarrows F_2
\]

$\pi_0(F_*, x) = $ ens.\ des \uncertain{classes} de lacets en $x$

\note{un lacet est dessiné à côté, une boucle refermée sur le point $x$}

\[
\widetilde{D}_{n'} \longrightarrow \widetilde{D}_n
\qquad\qquad
\begin{array}{c} x'' \leq x' \\ y'' \leq y' \end{array}
\ \in \widetilde{D}_{\ill{}}
\]

\section{Modèles, types, et ensembles ordonnés M-épinglés}

\note{ce titre est de l'édition ; le feuillet n'en porte pas. Les pages 6 à 16
sont la rédaction au propre, à l'encre, sous sa pagination à lui, 1 à 6, qui
figure en tête des pages 6, 8, 10, 12, 14 et 16 seulement — les rectos. Les
pages 7, 9, 11, 13 et 15 sont les versos de ces mêmes feuillets, ne portent
pas de numéro, et continuent directement le recto : le texte est donc continu
de la page 6 à la page 16}

\page{6}\note{sa page 1}

Soit $M$ une catégorie (« modèles »). On suppose

\begin{itemize}
\item[(i)] Toute flèche de $M$ est monomorphique
\item[(ii)] Tout \struck{automorphisme} \add{end.} d'un élément de $M$ est
l'identité.
\end{itemize}

Ex : $M$ une catégorie définie par un ensemble ordonné

Pour tout $D \in \mathrm{Ob}\,M$, soit $\widetilde{D}$ l'ensemble (ordonné par
« inclusion ») des sous-objets de $D$. On identifie $D$ au plus grand élément
de $\widetilde{D}$. Soient $D'$, $D \in \mathrm{Ob}\,M$, alors l'application

\[
(*) \qquad \mathrm{Hom}_M(D', D) \longrightarrow \widetilde{D},
\qquad u \longmapsto u(D')
\]

est injective par (i)(ii), son image est formée de l'ens.\ \struck{\ill{}}
$\Phi_{D'}(D) \subset \widetilde{D}$ des sous-objets $\Delta$ de $D$ qui sont
$M$-isomorphes à $D'$.

\struck{Soit $K$ un ens.\ ordonné}

Soit $\mathbb{N}$ l'ens.\ (préordonné) des classes d'isom.\ d'objets de $M$
\note{il est en fait ordonné, grâce à (ii) — l'incise est interlinéaire}
Pour $\forall\, n \in \mathbb{N}$, un \struck{s} objet $D_n$ de
\struck{$M$} de classe $n$ est déterminé à isom.\ unique près,
\uncertain{$\therefore$ comme dans (ii)} donc l'ens.\ ordonné
$\widetilde{D}_n = I_n$ est défini à isom.\ unique près. On a donc
\struck{\ill{}} une famille $(\widetilde{D}_n)_{n \in \mathbb{N}}$ d'ens.\
ordonnés, indexés par l'ensemble ordonné $\mathbb{N}$.

\page{7}

De plus, on se \struck{\ill{}} \uncertain{rappelle} la définition de

\[
I_n = \widetilde{D}_n = \text{ens.\ des sous-objets de } D_n
\]

--- et, pour $\forall\, n \in \mathbb{N}$, une application croissante,
strictement croissante

\[
I_n \overset{\tau_n}{\longrightarrow} \mathbb{N}
\qquad\qquad (\tau_n, \text{ application « type »})
\]

dont l'image est fournie \struck{dans} de l'ens.\
\struck{$I_{\leq n}$} \add{des} $\mathbb{N}_{\leq n} = \{n' \in \mathbb{N}
\mid n' \leq n\}$.

Notons que \uncertain{déjà} $I_n$ a un plus grand élément, qui est
« $D_n$ » (au sens \uncertain{que c'est} un sous-objet maximal de $D_n$), et
que je note $d_n$. On a donc (par (ii))

\[
\tau_n^{-1}(n) = \{d_n\}
\]

\struck{Si $n' \leq n$, alors pour tout $x \in I_n$}

Soient $n, n' \in \mathbb{N}$, \struck{$x \in I_n$, $n' = \tau_n(x)$,} avec
$n' \leq n$, on \struck{considérons} \struck{\ill{}} \add{considérons}
l'application

\[
\mathrm{Hom}(D_{n'}, D_n) \hookrightarrow \widetilde{D}_n = I_n,
\qquad u \longmapsto u(d_{n'})
\]

qui est une bijection des Hom sur l'ens.\ $\tau_n^{-1}(n')$, et considérons

\begin{tikzcd}[column sep=small, row sep=small, nodes={font=\scriptsize}]
\mathrm{Hom}(D_{n'}, D_n) \arrow[r] & \mathrm{Hom}(\widetilde{D}_{n'}, \widetilde{D}_n)
\end{tikzcd}

\note{sur le feuillet, $I_{n'}$ et $I_n$ sont écrits sous $\widetilde{D}_{n'}$
et $\widetilde{D}_n$}

on en déduit une \struck{bijec\ill{}} application

\[
\tau_n^{-1}(n') \overset{\text{déf}}{=} I_{n', \struck{n}}
\overset{\varphi_{n,n'}}{\longrightarrow} \mathrm{Hom}(I_{n'}, I_n),
\qquad \xi \longmapsto u_\xi^n
\]

\note{sous le membre de droite le feuillet porte
$[\mathrm{Hom}(D_{n'}, D_n)]$ surmonté d'un signe $\simeq$ ; à droite,
« une des éléments de $I_n$ de type $n'$ »}

qui est injective, \add{ou} plus précisément on a la commutativité de

\begin{tikzcd}[column sep=small, row sep=small, nodes={font=\scriptsize}]
I_{n,n'} \arrow[r, "\varphi_{n,n'}"] \arrow[rr, bend right=25, "\text{inclusion}"'] & \mathrm{Hom}(I_{n'}, I_n) \arrow[r, "u \mapsto u(d_{n'})"] & I_n
\end{tikzcd}

\page{8}\note{sa page 2}

i.e.\ on a, pour $\xi \in I_n(n')$ \add{i.e.\ $x \in I_n$, $\tau_n(x) = n'$}

\[
u_\xi(d_{n'}) = \xi.
\]

De plus, \ill{} on a : $u_\xi$ induit un \struck{\ill{}} isomorphisme
$I_{n'} \overset{\sim}{\longrightarrow} (I_n)_{\leq \xi}$.

Lorsque $n' = n$, on a $I_{n,n} = \{d_n\}$, et

\[
u_{d_n} = \mathrm{id}_{I_n}
\]

La connaissance de $\mathbb{N}$, $(I_n)_{n \in \mathbb{N}}$,
$(\tau_n)_{n \in \mathbb{N}}$, $(\varphi_{n,n'})_{(n,n') \in \mathbb{N} \times
\mathbb{N},\ n' \leq n}$ permet de reconstituer \struck{\ill{}} la catégorie
$M$, à \struck{l'équi} équivalence près ; il manque encore une
\struck{\ill{}} \add{condition} pour comprendre les homs.\ « standard » entre
les $I_n$, i.e.\ ceux provenant des $I_{n',n}$. \struck{Donc} On considère
leurs compositions

\begin{tikzcd}[column sep=small, row sep=small, nodes={font=\scriptsize}]
I_{n''} \arrow[r, "u_y^{n'}"] & I_{n'} \arrow[r, "u_\xi^n"] & I_n
\end{tikzcd}

\note{sur le feuillet, $y$ et $\xi$ sont écrits au-dessus de $I_{n'}$ et de
$I_n$ avec le signe d'appartenance, et la ligne $n'' \leq n' \leq n$ court
sous le diagramme}

on doit avoir

\[
u_\xi^n\, u_y^{n'} = u_\zeta^n
\qquad \text{où } \zeta = u_\xi^n(y) = (u_\xi^n\, u_y^{n'})(d_{n''})
\]

qui exprime la \struck{compatibilité} stabilité de la notion d'homs.\
admissibles pour les $I_n$ par composition.

\struck{Il est évident que la cat.} Soit $M_0$ la catégorie définie par

\[
\mathrm{Ob}\,M_0 = \mathbb{N},
\qquad
\mathrm{Hom}_{M_0}(n', n) =
\begin{cases}
\emptyset & \text{si } n' \nleq n \\
\text{ens.\ des homs.\ admissibles } I_{n'} \to I_n & \text{si } n' \leq n
\end{cases}
\]

\page{9}\note{sa page 3}

la composition des \add{admiss.} homs.\ étant celle des homs.\ entre les $I_n$
(qui sont des immersions fermées). On a un foncteur canonique (i.e.\ défini à
\add{un} isom.\ \ill{} près, $\xi_{n'} \longrightarrow \ill{}$)
\struck{(défini \ill{})}

\[
M \longrightarrow M_0
\]

qui est une équivalence de catégories --- en fait, c'est un isom.\ de
catégories, \struck{\ill{}} $M$ est une catégorie \uncertain{rigide} (deux
objets isom.\ sont identiques).

Je vais supposer par la suite $M$ rigide --- \uncertain{peut donc}
l'identifier : la catégorie $M_0$, définie par les données

\begin{itemize}
\item $\mathbb{N}$, ens.\ ordonné ;
\item $(I_n)_{n \in \mathbb{N}}$, famille d'ens.\ ordonnés, chaque $I_n$ ayant
un plus grand élément ;
\item $(\tau_n : I_n \to \mathbb{N})_{n \in \mathbb{N}}$, famille
d'applications strictement croissantes, $\tau_n(I_n) = \mathbb{N}_{\leq n}$,
$\tau_n^{-1}(n) = d_n$ (plus grand él.\ de $I_n$) ;
\item $(\varphi_{n,n'})_{(n',n) \in \mathbb{N} \times \mathbb{N},\ n' \leq n}$,
famille d'applications
\[
\varphi_{n,n'} : \tau_n^{-1}(n') \longrightarrow \mathrm{Hom}(I_{n'}, I_n),
\qquad \xi \longmapsto \bigl( u_\xi^n : I_{n'} \to I_n \bigr)
\]
satisfaisant
\[
u_\xi^n : I_{n'} \overset{\sim}{\longrightarrow} (I_n)_{\leq \xi}
\qquad (\text{d'où } u_\xi^n(d_{n'}) = \xi)
\]
\[
u_\xi^n \circ u_y^{n'} = u_\zeta^n
\quad \text{où } \zeta = u_\xi^n(y),
\quad \xi \in I_n,\ y \in I_{n'},\ n' = \tau_n(\xi),\ n'' = \tau_{n'}(y).
\]
\end{itemize}

\note{sur le feuillet ces quatre données sont rangées côte à côte sous une
grande accolade, chacune avec sa glose en dessous}

\page{10}\note{sa page 3, suite ; le numéro est écrit une seule fois en tête
du recto}

Inversement, de telles données définissent une catégorie $M_0$, satisfaisant à
(i), (ii).

N.B. Dans les cas que j'ai en tête, $\mathbb{N} = \mathbf{N}$.

\textbf{Ex 1} Catégorie des ens.\ finis tot.\ ordonnés $\Delta_n$, i.e.\ des
$\Delta_n = \ill{}$ et des applications strictement croissantes
\note{au-dessus de la ligne, souligné : « $n$-places types »}

\marginal{$I_n = \mathfrak{P}^*(\Delta_n)$, \quad Card $I_n = 2^{n+1} - 1$
\qquad ($n$-places types)}

\textbf{Ex 2} Catégorie des \struck{cubes} \struck{\ill{}} épinglés, i.e.\ des
\ill{} \struck{\ill{}} combinatoires, \struck{$A$, $\widetilde{A}$, \ill{},
$\widetilde{A} = A$} épinglés par un ordre total sur l'ens.\ des paires
d'hyperfaces, et le choix d'une \ill{} de chaque paire. \ill{} L'épinglage
d'un cube induit un épinglage de toute face, dans \ill{}

\marginal{$I_n = $ ens.\ des \ill{} $\Delta_{n-1} \times \{\pm 1\}$ \ill{}
d'ordres \ill{} épinglés \ill{} \quad Card $I_n = 3^n$}

\textbf{Ex. 3} Catégorie \struck{des} \add{des} \uncertain{hémisphères} $D_n$,
où

\[
I_n = (\Delta_n \times \{\pm 1\}) \sqcup \{d_n\},
\]

avec $\{\pm 1\}$ muni de l'ordre chaotique, $\{d_n\}$ le plus grand élément de
$I_n$.

\marginal{Card $I_n = 2n+1$}

\note{la marge gauche de ce feuillet est écrite en biais, à quelque quarante
degrés ; seuls les trois dénombrements reproduits ci-dessus et le membre
$\Delta_{n-1} \times \{\pm 1\}$ y portent. Ce dernier ne s'accorde pas avec le
$(\Delta_n \times \{\pm 1\}) \sqcup \{d_n\}$ du corps, dont le cardinal est
$2n+3$ et non $2n+1$ ; l'écart est signalé et non réparé}

Dans \struck{tous} \add{trois} ces exemples, on a
$I_n^* = I_n \smallsetminus \{d_n\}$

\begin{itemize}
\item[\textbf{a}] \add{si $n \geq 1$,} la réalisation top.\ $|I_n^*|$ de
$I_n^*$ \struck{\ill{}}
\item[\textbf{b}] \struck{\ill{}} est une $(n-1)$-sphère, (donc $|I_n|$ est un
$n$-disque)
\end{itemize}

\note{les deux items sont cerclés sur le feuillet et l'ouverture de chacun est
biffée : la phrase court d'un item à l'autre}

\page{11}

\begin{itemize}
\item[\textbf{b}] Les ens.\ \add{ordonnés} $I_n^*$, (regardés comme un \ill{}
polyèdre), sont \uncertain{réguliers}, i.e.\ le \ill{} est transitif sur
l'ens.\ des repères de $I_n^*$, i.e.\ l'ens.\ des chaînes maximales des
$(\xi_0 \subset \xi_1 \subset \cdots \subset \xi_{n-1})$
\end{itemize}

\textbf{Question} \quad Y a-t-il d'autres exemples d'une suite
\struck{$(P_n \subset I_n^*)$} $(P_n)_{n \geq 1}$
\struck{$(I_n^*)_{n \in \mathbb{N},\ n \geq 1}$} d'ens.\ ordonnés
\note{sous la ligne : $P_{n-1}$} qui \ill{} les propriétés

\begin{itemize}
\item[a)] $P_n$ est un $n$-polyèdre \uncertain{régulier} sphérique
\item[b)] Si $\xi \in P_n$, \struck{et de dimension} \add{alors}
$(P_n)_{\leq \xi}$ \struck{\ill{}} \struck{polyèdre combinatoire de dimension
$n'$, \ill{}}
\[
(P_n)_{\leq \xi} \simeq P_{n'}, \qquad \text{pour } n' \text{ \uncertain{convenable} } \leq n.
\]
\end{itemize}

(Ça doit être dans les livres de Coxeter !)

On aimerait de plus, surtout, être capables de \struck{\ill{}} des
\struck{\ill{}} $\varphi_{n,n'}$ satisfaisant à la condition de transitivité.

\marginal{b) est-il canonique ?! \quad c) \ill{} \qquad $P_n = I_{n+1}^*$}

\page{12}\note{sa page 4}

Je considère maintenant la sous-catégorie $\mathbb{B}_0$ de
$\widehat{M} \supset M$, formée des objets $X$ de la forme suivante :

\begin{itemize}
\item $X$ est \struck{réunion} \add{sup.} de ses sous-objets isomorphes à des
objets de $M$ (modèles)
\item \struck{et} si $K, L \subset X$ sont des sous-objets modèles, alors
$K \cap L$ est sup.\ de sous-objets de $K \cap L$ qui sont des modèles
\end{itemize}

Soit donc $\widetilde{X}$ l'ens.\ des \struck{\ill{}} sous-objets de $X$ qui
sont des modèles. Alors $\widetilde{X}$ \add{$K$} est un ens.\ ordonné, muni
de la structure suivante (« $M$-épinglage » de $K$) :

\begin{itemize}
\item[a)] \struck{\ill{}} une application strictement croissante
\[
K \overset{\tau_K}{\longrightarrow} \struck{\mathbb{N}}\ \mathbf{N}
\qquad (\text{application « type »})
\]
\item[b)] pour $\forall\, \xi \in K$, \add{pour $n = \tau_K(\xi)$}
\struck{une application} \struck{\ill{}}
\[
\struck{\varphi_\xi}\ I_n \overset{u_\xi^K}{\underset{\sim}{\longrightarrow}} K_{\leq \xi}
\qquad (\text{application d'épinglage de } K_{\leq \xi})
\]
\[
u_\xi^K(d_n) = \xi \quad (= \text{plus grand élément de } K_{\leq \xi})
\]
\end{itemize}

ces \struck{conditions} propriétés \struck{\ill{}} :

\marginal{les applications $I_n \to K$ définies par les $\tau^{-1}(n)$ et
dites \ill{} $(= K(n))$ \uncertain{admissibles}}

\page{13}

\[
(1) \qquad \tau_K \circ u_\xi^K = \tau_n
\]

et la propriété de transitivité pour les applications admissibles

\[
(2) \qquad \forall\, \xi \in K,\ y \in I_n \ (n = \tau_K(\xi)),
\quad \text{on a, si } n' = \tau_n(y),
\]
\[
\struck{H}\ u_\xi^K \circ u_y^n = u_\zeta^K,
\qquad \text{\add{où} } \zeta = u_\xi^K(y),
\]

mais \struck{(2) Si $\xi \in K$, $y \in K$, alors}

La catégorie dont les objets sont les objets admissibles de $\widehat{M}$, les
flèches les \uncertain{homs.} entre eux, est équivalente à \ill{} : celle des
ens.\ ordonnés $M$-épinglés.

D'autre part les objets quelconques de $\widehat{M}$ et leurs \ill{} se
reconstituent également, ce sont les \uncertain{applications strictement
croissantes}

\[
K \overset{f}{\longrightarrow} L
\]

telles que

\begin{itemize}
\item[a)] $\tau_L \circ f = \tau_K$
\item[b)] pour $\forall\, \xi \in K$, (on ait, \add{pour $y = f(\xi) \in L$})
la commutativité de
\end{itemize}

\begin{tikzcd}[column sep=small, row sep=small, nodes={font=\scriptsize}]
K_{\leq \xi} \arrow[rr, "f|K_{\leq \xi}"] & & K_{\leq y} \\
& I_n \arrow[ul, "\sim"] \arrow[ur, "\sim"'] &
\end{tikzcd}

\note{le sommet de droite est tracé comme un $K$ ; $f$ allant de $K$ dans $L$,
on attend $L_{\leq y}$, qui est ce qu'écrit la page 14. Les deux obliques ne
portent qu'un signe d'isomorphie ; ce sont $u_\xi^K$ et $u_y^L$. Où
$n = \tau_K(\xi) = \tau_L(y)$}

\page{14}\note{sa page 5}

ce qui implique en particulier que $f|K_{\leq \xi}$ induit un iso

\[
K_{\leq \xi} \overset{\sim}{\longrightarrow} L_{\leq y}
\]

Nous nous intéressons surtout au cas où $f$ est injectif, ce qui correspond
\struck{\ill{}} au cas \struck{des} des \uncertain{monos} dans $\widehat{C}$.

Alors \struck{\ill{}} $f$ induit \struck{\ill{}} un \uncertain{isom.} de $K$
sur \struck{\ill{}} son image $L' = f(K)$, muni de l'ordre induit par $L$.
Notons que $L'$ est une partie fermée de $L$ ($x \in L'$, $y \leq x
\Rightarrow y \in L'$) et que les parties fermées admettent un $M$-épinglage
induit \struck{\ill{}} unique, tel que l'inclusion soit admissible (épinglage
induit) --- et remarquons\add{:} que $f$, bien sûr
$K \overset{f}{\underset{\sim}{\longrightarrow}} L'$, est compatible avec les
$M$-épinglages.

\page{15}

Un cas intéressant est celui où $K$ \struck{\ill{} est la partie formée par
\ill{}} et \ill{} objets modèles, ce qui signifie que $K$ \ill{} est fini, ou
du moins \add{n'a qu'un nombre fini d'}éléments maximaux, et que $K$ est
engendré par ses éléments maximaux. Pour un tel ens.\ \ill{} \add{le
\ill{}} \ill{} qui a cette propriété, un $M$-épinglage est connu quand on
connaît les $u_\alpha^K$ pour $\alpha \in \ill{}$ des \ill{}.

La condition pour que les

\[
u_\alpha^K : I_{\tau_n(\alpha)} \overset{\sim}{\longrightarrow} K_{\leq \alpha}
\qquad (\alpha \in \mathrm{Max}\,K)
\]

définissent un épinglage, est que pour \struck{\ill{}} $\alpha, \beta \in
\mathrm{Max}\,K$ et $\xi \in K$, $\xi \leq \alpha, \beta$, posant
$n = \tau_K(\xi)$, les deux iso

\[
I_n \overset{\sim}{\longrightarrow} K_{\leq \xi}
\]

définis par l'épinglage \struck{\ill{}} $u_\alpha^K$ de $K_{\leq \alpha}$, et
par l'épinglage $u_\beta^K$ de $K_{\leq \beta}$, soient les
\uncertain{mêmes}.

\page{16}\note{sa page 6}

La catégorie des objets admissibles de $\widehat{C}$ est stable par
\struck{les} \struck{l'opération} sommes quelconques, et par sommes
amalgamées $X \sqcup_Z Y$, dans \ill{} $Z \to X$ et $Z \to Y$ des monos, et
\add{de la} le foncteur $X \longrightarrow \widetilde{X}$ commutant à ces
opérations, \ill{} de façon évidente. De plus, les objets de type fini
donnent des objets de t.f., pourvu, toutefois, \add{dans} le cas des sommes,
qu'elles se bornent à des sommes \uncertain{finies}.

N.B. Quand les $I_n$ sont finis, alors les objets \add{adm.} de t.f.\ de
$\widehat{C}$ \struck{\ill{}} correspondent aux ens.\ ordonnés
$M$-épinglés $K$ qui sont \uncertain{finis}.

\end{document}
